\section{结论}

本项目完成了一个基于Qiskit、Streamlit和Plotly的Grover搜索算法振幅放大可视化系统。系统支持2到3个量子比特，能够选择任意目标态，并在Hadamard初始化、Oracle相位翻转和Diffusion均值倒置后保存并展示中间statevector。

通过\(n=3\)、目标态\(|101\rangle\)的实验可以看到，Oracle只是将目标态概率幅变为负值，目标态概率并未立即增加；随后Diffusion通过关于平均概率幅的反演，把相位标记转化为振幅放大，使目标态概率从\(12.5\%\)逐步提升到\(78.125\%\)和\(94.53125\%\)。这一结果与Grover算法理论公式一致。

因此，本项目较好地完成了课程选题要求：利用柱状图和3D步进图展示Grover搜索算法中目标态概率幅被逐步放大的过程，并通过交互界面降低了抽象量子态演化的理解难度。
